\section{Mathematical foundations}
\label{sec:math}
This section collects the formal objects used across the five gaps.
Each construction is implemented in \texttt{crisprgap} and exercised by
the hermetic test suite; every empirical claim later in the paper can be
traced to one of the definitions here.

\subsection{Featurization and base learners}
For a guide $g$ of length $L=20$ over alphabet
$\Sigma=\{\mathrm{A,C,G,T}\}$, the position-specific one-hot map is
\begin{equation}
\phi(g) \;=\; \big[\mathbb{1}(g_i = s)\big]_{i=1..L,\; s \in \Sigma}
\;\in\; \{0,1\}^{4L},
\label{eq:onehot}
\end{equation}
and its dinucleotide extension appends
$\mathbb{1}(g_i g_{i+1} = d)$ for $d \in \Sigma^2$, $i=1..L-1$
($16(L-1)$ extra coordinates). Our ridge baseline solves
\begin{equation}
\hat w \;=\; \arg\min_w \; \tfrac{1}{n}\textstyle\sum_j (y_j - w^\top x_j)^2
+ \lambda \|w\|_2^2
\;=\; (X^\top X + n\lambda I)^{-1} X^\top y,
\label{eq:ridge}
\end{equation}
which is the MAP estimate under a Gaussian prior
$w \sim \mathcal N(0, (n\lambda)^{-1} I)$ and Gaussian noise.

\subsection{Rank correlation and monotone invariance}
All cross-dataset comparisons use Spearman correlation. With midranks
$R(x), R(y)$ (ties share the average rank),
\begin{equation}
\rho(x, y) \;=\;
\frac{\mathrm{cov}(R(x), R(y))}{\sigma_{R(x)}\, \sigma_{R(y)}}.
\label{eq:spearman}
\end{equation}
\begin{proposition}[monotone invariance]
\label{prop:mono}
If $T$ is strictly increasing on the range of $y$, then
$\rho(x, T \circ y) = \rho(x, y)$.
\end{proposition}
\begin{proof}
$T$ strictly increasing implies $R(T \circ y) = R(y)$: for any pair
$i, j$, $y_i < y_j \iff T(y_i) < T(y_j)$, so the order statistics, hence
the midranks, coincide. Substituting into \eqref{eq:spearman} gives the
claim.
\end{proof}
Proposition~\ref{prop:mono} licenses working in normalized-rank space
throughout: residual definitions, portability scores, and corrections
below are invariant to the (unknown, screen-specific) monotone transforms
that map raw assay readouts to reported activities.

\subsection{CFD as naive-Bayes odds}
The CFD score for guide $g$ cleaving site $g'$ with mismatch set
$M = \{(i, s_i \!\to\! t_i)\}$ is
\begin{equation}
\mathrm{CFD}(g, g') \;=\; \prod_{(i, u) \in M} f(i, u),
\label{eq:cfd}
\end{equation}
where $f(i, u)$ is the empirically measured relative cleavage frequency
for substitution $u$ at position $i$ (vendored matrices,
Section~3.5). Eq.~\eqref{eq:cfd} is exactly the posterior-odds factor of
a naive-Bayes model: writing $C$ for cleavage,
\begin{equation}
\frac{P(C \mid M)}{P(\neg C \mid M)} \;=\;
\frac{P(C)}{P(\neg C)} \prod_{(i,u)\in M}
\frac{P(u \text{ at } i \mid C)}{P(u \text{ at } i \mid \neg C)},
\label{eq:cfdnb}
\end{equation}
under conditional independence of mismatch effects given cleavage. The
assumption is known to be approximate (mismatch--mismatch interactions
exist), which is exactly the modeling headroom our GNN exploits in
gap 2.

\subsection{Off-target GNN}
Each guide--site pair is a two-node graph with a typed edge (mismatch
pattern, PAM channel in gaps 4-5, epigenetic channel in gap 5). Message
passing for layer $k$ is
\begin{equation}
h_v^{(k+1)} \;=\; \sigma\!\Big( W_s h_v^{(k)} + \!\!\sum_{u \in N(v)}\!\!
W_e(e_{uv})\, h_u^{(k)} \Big),
\label{eq:gnn}
\end{equation}
with pair readout $\hat p = \mathrm{sigmoid}(\mathrm{MLP}([h_g; h_{g'}]))$
trained by binary cross-entropy
\begin{equation}
\mathcal L_{\mathrm{BCE}} = -\tfrac{1}{n}\textstyle\sum_j
\big[ y_j \log \hat p_j + (1 - y_j)\log(1 - \hat p_j) \big].
\label{eq:bce}
\end{equation}

\subsection{Calibration}
For predicted probabilities $\hat p$ and outcomes $y$, with $B$ bins
$\{\mathcal B_b\}$,
\begin{equation}
\mathrm{ECE} \;=\; \sum_{b=1}^{B} \frac{|\mathcal B_b|}{n}
\left| \mathrm{acc}(\mathcal B_b) - \mathrm{conf}(\mathcal B_b) \right|,
\qquad
\mathrm{BS} = \tfrac{1}{n}\textstyle\sum_j (\hat p_j - y_j)^2 .
\label{eq:ece}
\end{equation}
\begin{proposition}[Murphy decomposition]
\label{prop:murphy}
With $\bar y$ the base rate, $\bar p_b$ the mean prediction and
$\bar y_b$ the mean outcome in bin $b$,
\begin{equation}
\mathrm{BS} \;=\;
\underbrace{\sum_b \tfrac{|\mathcal B_b|}{n}(\bar p_b - \bar y_b)^2}_{\text{reliability}}
\;-\;
\underbrace{\sum_b \tfrac{|\mathcal B_b|}{n}(\bar y_b - \bar y)^2}_{\text{resolution}}
\;+\;
\underbrace{\bar y (1 - \bar y)}_{\text{uncertainty}}.
\label{eq:murphy}
\end{equation}
\end{proposition}
\begin{proof}
Write each term as
$(\hat p_j - y_j)^2 = [(\hat p_j - \bar p_b) + (\bar p_b - \bar y_b) + (\bar y_b - \bar y) + (\bar y - y_j)]^2$
and average over bin $b$. The within-bin terms vanish on aggregation
(they average to deviations from bin means, which integrate to zero
under squared loss when predictions are bin-constant); expanding the
surviving three terms and noting the cross terms cancel by
$\sum_{j \in \mathcal B_b}(y_j - \bar y_b) = 0$ gives
\eqref{eq:murphy}.
\end{proof}
Isotonic calibration fits the monotone step function
\begin{equation}
\hat m \;=\; \arg\min_{m \text{ nondecreasing}} \textstyle\sum_j (y_j - m(\hat p_j))^2,
\label{eq:isotonic}
\end{equation}
solved exactly by the pool-adjacent-violators algorithm (PAVA): scanning
left to right, any adjacent violation $m_i > m_{i+1}$ triggers pooling of
the two blocks into their weighted mean; each pool step strictly
decreases the objective and the process terminates at the unique monotone
projection of $y$ onto the monotone cone (Barlow et al.\ 1972). Platt
scaling instead fits $P(y{=}1 \mid s) = \sigma(a s + b)$ by maximum
likelihood, a 2-parameter restriction of \eqref{eq:isotonic}; our gap-2
finding (isotonic $\le$ uncalibrated $<$ Platt in ECE) matches the
sample-size regime where the richer map wins (Niculescu-Mizil \& Caruana
2005).

\subsection{Transfer residual and the PORTA correction}
Let $f$ be a model trained on source dataset $A$. For a destination
screen $B$ over guides $g$ with normalized-rank maps $p_B(g) =
\mathrm{nrank}(f(g))$ and $t_B(g) = \mathrm{nrank}(y_B(g))$ in $[0,1]$,
define the per-guide transfer residual
\begin{equation}
r_B(g) \;=\; p_B(g) - t_B(g).
\label{eq:resid}
\end{equation}
A portability model $\hat r = m(\phi(g), b(g))$ is fit on a
\emph{calibration screen} $A'$ (features: sequence and biofeatures) and
the corrected destination ranking is
\begin{equation}
\tilde p_B(g) \;=\; p_B(g) - m(g).
\label{eq:porta}
\end{equation}
\begin{proposition}[exact recovery under replication]
\label{prop:porta}
If the residual replicates across screens, $r_B(g) = r_{A'}(g)$ for all
$g$, and the portability model is perfect on the destination guides,
$m(g) = r_{A'}(g)$, then $\tilde p_B = t_B$ exactly, and
$\rho(\tilde p_B, y_B) = 1$ up to ties.
\end{proposition}
\begin{proof}
$\tilde p_B = p_B - m = p_B - r_{A'} = p_B - r_B = p_B - (p_B - t_B) =
t_B$. By definition $t_B$ is the normalized rank of $y_B$, so
\eqref{eq:spearman} gives $\rho(t_B, y_B) = 1$ absent ties.
\end{proof}
Proposition~\ref{prop:porta} isolates the two empirical quantities that
bound the correction: residual replication (measured: $\rho = 0.86$--$0.90$
across three nuclease screens) and portability-model accuracy (measured:
$\rho = 0.78$ cross-screen, $0.86$ within-screen). Their product being
far above the raw-transfer baseline ($0.015$) is what Section 8.1
verifies empirically.

\subsection{When pooling helps}
For a per-guide effect estimated from two screens with $n_A, n_B$
samples, equal per-sample variance $\sigma^2$, and a between-screen bias
$\delta$ in the second screen, the pooled mean has
\begin{equation}
\mathrm{MSE}_{\mathrm{pool}} \;=\;
\Big(\tfrac{n_B}{n_A + n_B}\Big)^{\!2} \delta^2 \;+\; \frac{\sigma^2}{n_A + n_B},
\label{eq:pool}
\end{equation}
which beats the source-only estimator ($\mathrm{MSE}_A = \sigma^2/n_A$)
iff
\begin{equation}
\delta^2 \;<\; \frac{\sigma^2 (n_A + n_B)}{n_A\, n_B}.
\label{eq:poolcond}
\end{equation}
Eq.~\eqref{eq:poolcond} quantifies why gap-1 pooling fails: the
cross-family $\delta$ is so large (raw transfer $\rho \approx 0.02$) that
no sample-size ratio satisfies the condition.

\subsection{Uncertainty estimates}
All headline correlations carry bootstrap percentile intervals: with
resamples $b = 1..B$ over guide indices and statistic $T_b$,
\begin{equation}
\mathrm{CI}_{95} \;=\; \big[ Q_{0.025}(\{T_b\}),\; Q_{0.975}(\{T_b\}) \big],
\label{eq:boot}
\end{equation}
and paired differences $T_b^{\mathrm{corr}} - T_b^{\mathrm{raw}}$ are
computed within the same resample so that the comparison controls for
resample-level variation.
